% PAP-RES-001  (JMFK)
% Variable projection fitting of dark adaptation recovery curves:
% method and numerical validation.
% Compile with: xelatex -> bibtex -> xelatex -> xelatex
% Requires varpro.bib alongside this file.
\documentclass[11pt,a4paper]{article}
\usepackage[margin=25mm]{geometry}
\usepackage{fontspec}
\IfFontExistsTF{Helvetica Neue}
  {\setmainfont{Helvetica Neue}}
  {\relax} % fall back to the default font if Helvetica Neue is absent
\usepackage{amsmath}
\usepackage{booktabs}
\usepackage{fancyhdr}
\usepackage{graphicx}
\usepackage[framemethod=tikz]{mdframed}
\usepackage{enumitem}
\usepackage{natbib}
\bibliographystyle{unsrtnat}
\usepackage[hidelinks]{hyperref}

% ---- colours, matching DISC-IP-003 ----
\definecolor{coneC}{HTML}{D85A30}
\definecolor{rodC}{HTML}{534AB7}
\definecolor{tealC}{HTML}{0F6E56}
\definecolor{grayC}{HTML}{5F5E5A}

% ---- document identity ----
\newcommand{\docid}{PAP-RES-001}
\newcommand{\docversion}{v0.1 skeleton}
\newcommand{\docdate}{27 August 2026}

% ---- headers and footers ----
\pagestyle{fancy}
\fancyhf{}
\fancyhead[L]{\small MuMac Ltd}
\fancyhead[R]{\small \docid\ \docversion}
\fancyfoot[L]{\footnotesize Draft manuscript. Not for circulation.}
\fancyfoot[R]{\small Page \thepage}
\renewcommand{\headrulewidth}{0.4pt}

% ---- placeholder boxes ----
\newmdenv[linecolor=coneC, linewidth=1pt, roundcorner=4pt,
  backgroundcolor=coneC!6, frametitle={To write},
  frametitlefont=\small\bfseries]{towrite}
\newmdenv[linecolor=tealC, linewidth=1pt, roundcorner=4pt,
  backgroundcolor=tealC!6, frametitle={Evidence in hand},
  frametitlefont=\small\bfseries]{evidence}
\newmdenv[linecolor=rodC, linewidth=1pt, roundcorner=4pt,
  backgroundcolor=rodC!6, frametitle={Open question},
  frametitlefont=\small\bfseries]{openq}

% ---- inline placeholder for a number not yet extracted ----
\newcommand{\FILL}[1]{\textbf{\textcolor{coneC}{[#1]}}}

% ---- model shorthand ----
\newcommand{\Sii}{S_2}

\begin{document}

% ================= title block =================
\begin{center}
{\LARGE Variable projection fitting of dark\\adaptation recovery curves}\\[3mm]
{\large Method, and numerical validation across\\192,000 simulated recoveries}\\[6mm]
J.\,M.\,F.~Kelly\\
\small MuMac Ltd, Manchester, United Kingdom\\[6mm]
\begin{tabular}{ll}
\toprule
Document ID & \docid \\
Version & \docversion \\
Date & \docdate \\
Author & JMFK \\
Status & Skeleton for drafting \\
\bottomrule
\end{tabular}
\end{center}

\vspace{3mm}
\noindent\textbf{Version history}

\noindent\begin{tabular}{llll}
\toprule
Version & Date & Author & Change \\
\midrule
v0.1 & \docdate & JMFK & Section skeleton and result placeholders \\
\bottomrule
\end{tabular}

\vspace{6mm}

% ================= abstract =================
\section*{Abstract}
\begin{towrite}
Write last. Structure: rod-mediated dark adaptation is an early functional
biomarker of AMD, but the incumbent endpoint (rod intercept time) is
censored whenever recovery outruns the test window, and the parameter-level
alternative has been considered unusable on grounds of fitting instability.
We show the recovery model is separable, give an exhaustive profiled fit
with bounded deterministic cost, and validate it over 192,000 simulated
recoveries. State the four headline results: gradient descent contributes
nothing measurable; the interval-based detection rule is robust where the
model-comparison rule fails; the detection threshold collapses to
$|\Sii|_{90} \approx 0.55\sigma$ per minute; and fits agree to $10^{-14}$
across two unrelated hardware platforms. Target 200 words.
\end{towrite}

% ================= 1 =================
\section{Introduction}

\subsection{Rod intercept time and its censoring problem}
\begin{towrite}
RIT as the incumbent clinical endpoint \citep{Jackson:2014aa,Owsley:2016aa}.
Why it is attractive: single number, monotone, clinically interpretable.
Why it is limited: it is an observation of when a curve crosses a criterion,
so a recovery slower than the test window yields no value at all. The field's
responses have been to extend the test, shorten it and predict
\citep{Jackson:2008aa,DeSilva:2019aa}, or to treat the missing values as
censored and apply survival methods \citep{Higgins:2018aa}.
\end{towrite}

\subsection{Why parameter-level estimation has been avoided}
\begin{towrite}
The rod slope $\Sii$ is the quantity the clinician actually wants: a rate,
not a crossing time, and defined whether or not the criterion is reached
within the window. The obstacle is that recovering it requires fitting a
segmented nonlinear model, and the received view is that such fits are
unstable, initialisation-dependent and not reproducible. That view is
correct about the optimisers usually applied and wrong about the model.
\end{towrite}

\subsection{Contribution}
\begin{enumerate}[nosep]
\item The recovery model is separable: with the two nonlinear parameters
  held, the remaining three solve in closed form.
\item An exhaustive search over the two-dimensional nonlinear space then
  gives a fit with bounded, predictable cost, no dependence on starting
  values, and no differentiation of the breakpoint.
\item A validation of 192,000 replications establishing what the method
  does and does not deliver, including the detection threshold as a single
  scaling law and a demonstration of bit-level cross-platform agreement.
\end{enumerate}

% ================= 2 =================
\section{The recovery model}

\subsection{Functional form}
The five-parameter cone-plus-rod model, referred to here as P5c, is
\begin{equation}
\label{eq:p5c}
y(t) \;=\; ct \;+\; cc\,e^{-t/\tau} \;+\; \Sii\,\max(t-\alpha,\,0),
\end{equation}
with $ct$ the cone threshold asymptote, $cc$ the cone coefficient, $\tau$
the cone time constant, $\Sii$ the rod slope and $\alpha$ the rod-cone break.

\begin{towrite}
Physiological reading of each parameter. State that $\Sii$ is the diagnostic
parameter and everything else is, for present purposes, nuisance structure
that must nevertheless be estimated well enough not to contaminate it.
\end{towrite}

\subsection{The transition function}
\begin{towrite}
Equation~\eqref{eq:p5c} uses a hinge. Earlier work, including the granted
patent \citep{Kelly:2020aa}, used a logistic transition $Tr(\cdot)$.
Justify the hinge: the smooth transition was introduced to make the model
differentiable for gradient-based fitting, and once the fit no longer
differentiates with respect to the breakpoint that motivation disappears.
State what is lost, and whether the two forms are distinguishable at
clinical noise. \FILL{decide whether this needs its own small experiment}
\end{towrite}

\subsection{The cone-only null model}
\begin{towrite}
P3: $y(t) = ct + cc\,e^{-t/\tau}$. Needed for the model-comparison detection
rule assessed in section~\ref{sec:rule}.
\end{towrite}

% ================= 3 =================
\section{Method}

\subsection{Separability}
Hold $\tau$ and $\alpha$ fixed. Equation~\eqref{eq:p5c} is then linear in
$ct$, $cc$ and $\Sii$, with design matrix
\begin{equation}
X(\tau,\alpha) \;=\;
\begin{bmatrix}
1 & e^{-t_1/\tau} & \max(t_1-\alpha,0)\\
\vdots & \vdots & \vdots\\
1 & e^{-t_n/\tau} & \max(t_n-\alpha,0)
\end{bmatrix},
\end{equation}
so the conditional least-squares solution is obtained directly by QR
decomposition and the objective reduces to a function of $(\tau,\alpha)$
alone.

\begin{towrite}
State this as variable projection \citep{Golub:1973aa,Golub:2003aa} and note
the specific consequence for this model: the parameter that the literature
treats as the hard one, the breakpoint, never needs to be differentiated.
\end{towrite}

\subsection{Exhaustive grid over the nonlinear pair}
\begin{towrite}
Grid specification, bounds and step sizes for $\tau$ and $\alpha$. Cost is
$n_\tau \times n_\alpha$ QR solves, known before the fit begins. Note the
relationship to grid search for segmented regression
\citep{Hudson:1966aa,Lerman:1980aa}, and what is new here: the segmented
term is one component of a model that is also nonlinear elsewhere.
\end{towrite}

\begin{table}[h]
\centering
\caption{Search grid and bounds. \FILL{fill from dc\_run\_experiment\_driver.R}}
\begin{tabular}{lrrr}
\toprule
Parameter & Lower & Upper & Step \\
\midrule
$\tau$ (min) & \FILL{} & \FILL{} & \FILL{} \\
$\alpha$ (min) & \FILL{} & \FILL{} & \FILL{} \\
\bottomrule
\end{tabular}
\end{table}

\subsection{Bound handling and the constrained re-solve}
\begin{towrite}
$\Sii$ is constrained to be non-positive on physiological grounds. When the
unconstrained conditional solution violates that, the exact constrained
solution is obtained by re-solving with the rod column dropped rather than
by clamping. State why clamping is wrong: it leaves $ct$ and $cc$ at values
optimal for a different model.
\end{towrite}

\subsection{Local refinement}
\begin{towrite}
Optional Nelder-Mead polish within the winning cell, and the gradient
descent stage that the validation subsequently shows to be redundant
(section~\ref{sec:gd}). Both are described here as implemented, and
adjudicated in results.
\end{towrite}

\subsection{Intervals}
\begin{towrite}
Delta-method intervals on $\Sii$ and $\alpha$, and the propagated interval
on predicted RIT. Note honestly that this remains conceptually close to the
approach of \citet{Kelly:2020aa}.
\end{towrite}

% ================= 4 =================
\section{Related work}
\begin{towrite}
Four threads, one paragraph each.
\begin{enumerate}[nosep]
\item Separable least squares \citep{Golub:1973aa,Golub:2003aa}.
\item Segmented regression by grid search
  \citep{Hudson:1966aa,Lerman:1980aa}, and the known asymmetry between
  coefficient and change-point interval coverage \citep{Kim:2008aa}, which
  section~\ref{sec:ci} reproduces.
\item Alternative recovery models, in particular the softplus form of
  \citet{Dimitrov:2008aa}, $T(t) = T_f + \log_{10}(10^{T_i - Rt} + 1)$,
  which is partially separable but not usefully so. \FILL{decide whether to
  fit it as a comparator or discuss only}
\item Clinical dark adaptometry and the RIT literature
  \citep{Jackson:2008aa,Jackson:2014aa,Owsley:2016aa,DeSilva:2019aa,Higgins:2018aa}.
\end{enumerate}
\end{towrite}

% ================= 5 =================
\section{Validation design}

\subsection{Generative model and factor levels}
\begin{towrite}
Curves simulated from equation~\eqref{eq:p5c} with additive Gaussian noise
on threshold. Fully crossed design over four cone time constants, six noise
levels and a range of rod slopes and break times, replicated to 192,000
runs in total.
\end{towrite}

\begin{table}[h]
\centering
\caption{Simulation factors. \FILL{fill exact level sets and replication count per cell}}
\begin{tabular}{lll}
\toprule
Factor & Levels & Note \\
\midrule
$\tau$ (min) & 0.6, 1.2, 1.8, 2.4 & 2.4 is a stress test; see section~\ref{sec:limits} \\
Noise $\sigma$ (log units) & \FILL{} & includes 0.05, 0.1, 0.3 \\
$\Sii$ (log units/min) & \FILL{} & spans shallow to steep \\
$\alpha$ (min) & \FILL{} & \\
Replications per cell & \FILL{} & 192,000 total \\
\bottomrule
\end{tabular}
\end{table}

\subsection{Reproducibility apparatus}
\begin{evidence}
Each result row carries its seed, the git commit of the fitting code, and a
captured environment record (platform, BLAS, R version, RNG version, git
state). The driver refuses to run against a dirty working tree. Output files
are named by PID, $\tau$ set and timestamp so concurrent runs cannot collide.
\end{evidence}

\subsection{Platforms}
\begin{evidence}
Two unrelated platforms: an Apple silicon Mac (arm64, Accelerate/vecLib
BLAS, R 4.5.2) and an UltraRISC UR-DP1000 RISC-V board (Ubuntu, reference
BLAS, R 4.3.3). \texttt{RNGversion("4.3.3")} was pinned on both and the
generated streams verified bit-identical via
\texttt{sprintf("\%a", rnorm(5))}, so any divergence in results is
attributable to arithmetic and not to the random numbers.
\end{evidence}

% ================= 6 =================
\section{Results}

\subsection{Gradient descent contributes nothing measurable}
\label{sec:gd}
\begin{evidence}
Median RSS improvement zero at all realistic noise levels. The optimiser
reports \texttt{no\_improving\_step} in 21 to 53\% of runs, consistent with
it being handed a point that is already stationary. Paired truth-error
comparison shows \FILL{alpha\_gain, S2\_gain medians and the fraction of
runs where GD is nearer truth; near 0.5 is a coin flip}.
\end{evidence}
\begin{towrite}
Frame as: the exhaustive grid is sufficient. This is the strongest form of
the robustness argument, because it removes the stage on which
initialisation dependence would otherwise rest.
\end{towrite}

\subsection{Nelder-Mead refinement is a real but small polish}
\begin{evidence}
Median RSS improvement 0.02 to 0.4\%. \FILL{effect on truth error, and the
$\tau$ and noise dependence}.
\end{evidence}

\subsection{Detection rule: intervals against model comparison}
\label{sec:rule}
\begin{evidence}
The AICc ceiling degrades monotonically in both noise and $\tau$: at noise
0.3 it falls 0.958, 0.887, 0.779, 0.698 across $\tau$ 0.6, 1.2, 1.8, 2.4,
with failures beginning inside the physiological range at $\tau$ 1.8. The
interval rule never falls below 0.988 across 24 cells, its 90\% thresholds
move by at most 0.015 log units/min over a fourfold range of $\tau$, and its
false-positive rate stays between 1 and 6\%. At $\tau$ 0.6 AICc is the better
detector (threshold 0.0125 against 0.022 at noise 0.05); by $\tau$ 1.8 it has
lost, and at $\tau$ 2.4 it does not reach 90\% anywhere above noise 0.1.
\end{evidence}
\begin{towrite}
The mechanism: a slow cone tail extends through the break region, the
exponential in P3 absorbs rod-branch curvature, and $\Delta$AICc starves.
The headline is the inversion. The model-comparison rule fails preferentially
on patients with slowed cone recovery, which is itself an early-AMD feature,
so it fails hardest on exactly the population the instrument exists to find.
That is why the interval rule ships, and the argument is robustness rather
than average performance.
\end{towrite}

\begin{table}[h]
\centering
\caption{Detection ceiling by rule, cone time constant and noise.}
\begin{tabular}{lrrrr}
\toprule
& \multicolumn{4}{c}{$\tau$ (min)} \\
\cmidrule(l){2-5}
Rule and noise & 0.6 & 1.2 & 1.8 & 2.4 \\
\midrule
AICc, $\sigma=0.3$ & 0.958 & 0.887 & 0.779 & 0.698 \\
Interval, $\sigma=0.3$ & \FILL{} & \FILL{} & \FILL{} & \FILL{} \\
\FILL{add remaining noise levels} & & & & \\
\bottomrule
\end{tabular}
\end{table}

\subsection{The detection threshold is a scaling law}
\begin{evidence}
Across all four cone sets, $|\Sii|_{90}/\sigma$ lies between 0.53 and 0.57
per minute at every noise level from 0.1 upward. The detectable slope is
approximately $0.55\sigma$ per minute and is essentially independent of
$\tau$.
\end{evidence}
\begin{towrite}
Convert to the clinical statement: at instrument noise of 0.1 log units, rod
slopes of 0.055 log units/min are detected nine times in ten at 3 to 5\%
false positives. Note that a twenty-cell table reduces to one number, and
that this is what makes the threshold quotable in a protocol.
\end{towrite}

\subsection{Interval calibration}
\label{sec:ci}
\begin{evidence}
Coverage of the $\Sii$ interval sits at \FILL{} against a nominal 0.95;
coverage of the $\alpha$ interval sits at \FILL{}. The asymmetry is the one
\citet{Kim:2008aa} report by simulation for segmented line regression.
\end{evidence}
\begin{towrite}
Present as a predicted result confirmed, not as a surprise. It validates the
interval machinery for the parameter the detection rule uses, and warns
against quoting a break-time interval at face value.
\end{towrite}

\subsection{Degenerate fits are made visible, not eliminated}
\begin{evidence}
8.0\% of fits terminate at the $\tau$ bound with $cc$ running to order
$10^{13}$, in simulated data with a known, well-defined $\tau$. The
degeneracy is a property of the model plus noise, not of the optimiser.
\FILL{breakdown by $\tau$ set and noise level}.
\end{evidence}
\begin{towrite}
The claim to make, and the one not to. What the exhaustive search delivers
is that the degenerate solution is known to be the global optimum rather
than a local trap, that the boundary condition is detected and flagged
deterministically instead of being returned silently as an excellent fit,
and that the same input always yields the same verdict. Do not write that
the method solves the $cc$/$\tau$ instability. It does not, and the data
above say so.
\end{towrite}

\subsection{Cross-platform agreement}
\begin{evidence}
Of 191,996 paired comparisons, 76.5\% of fits are bit-identical, 99\% agree
to $1\times10^{-14}$ in RSS, and the boundary classification disagrees in 1
comparison. All divergence traces to the Nelder-Mead accept test in cases
where the candidate improvement is at rounding level, and concentrates at
$\tau = 2.4$ (disagreement counts 1, 1, 1, 53 across the four sets).
\end{evidence}
\begin{towrite}
Why this matters for a regulated instrument: the fit is not merely
repeatable on one machine but reproducible across architectures and BLAS
implementations, which is the property a verification protocol needs. Note
that the one divergence is in the refinement stage, not the grid, so the
determinism claim attaches to the part of the method that does the work.
\end{towrite}

\subsection{RIT error and its propagation}
\begin{evidence}
Predicted RIT with propagated interval, compared against RIT computed from
the true parameters, by $\tau$ set. Medians are under a minute while 90th
percentiles reach 19 to 83 minutes: the $\Sii^{-4}$ tail. At
$\Sii = -0.05$ the propagated interval is approximately 5.5 min.
\end{evidence}
\begin{towrite}
The honest framing of the survival-analysis question. The method yields an
estimate with a stated interval where the observation would otherwise be
censored, which is an alternative to survival methods
\citep{Higgins:2018aa}. It does not eliminate the problem, it replaces
censoring with wide extrapolation, and it does not address their claim,
which is about statistical power in group comparisons. That is testable by
comparing required sample sizes, and should be offered as such.
\end{towrite}

\subsection{Sampling density}
\begin{evidence}
Thinning from 4/min to 2/min, 1/min and 0.5/min, method held constant.
Moving to 1/min, the AdaptDx density, costs roughly a factor of four in the
shallowest detectable slope, and halving again breaks detection at that
noise level. False positives stay between 0.3 and 2.5\% at every density and
noise level with no systematic trend. RIT error at 1/min is a median of
0.32 min at noise 0.1 and 0.89 min at noise 0.3.
\end{evidence}
\begin{towrite}
Two conclusions. First, the continuous ramp is not a convenience, it is
doing measurable diagnostic work, and the rule degrades by losing
sensitivity rather than by inventing rod phases, which is the right failure
direction for a screening instrument. Second, at matched density the RIT
error is two orders of magnitude below the 19.1~$\pm$~18.6 min of
\citet{DeSilva:2019aa}, so sampling density does not explain that
comparator's performance. See section~\ref{sec:limits} for the qualification
this needs.
\end{towrite}

% ================= 7 =================
\section{Discussion}
\begin{towrite}
Assemble in this order.
\begin{enumerate}[nosep]
\item The obstacle to parameter-level dark adaptometry was the fit, and the
  fit was harder than it needed to be because the model was treated as
  wholly nonlinear.
\item With separability exploited, the fit is deterministic, bounded in
  cost, free of starting values, and reproducible across architectures.
\item The parameters are therefore usable, so the sidestep to a crossing
  time is no longer forced.
\item The detection threshold is a single scaling law, which makes the
  instrument's noise specification and its clinical sensitivity two views of
  the same number.
\item The rule that ships is chosen for robustness on the population of
  clinical interest, not for average performance.
\end{enumerate}
Decide before drafting which sense of ``superior'' the paper commits to,
since that choice determines how much of this section exists.
\end{towrite}

% ================= 8 =================
\section{Limitations}
\label{sec:limits}
\begin{openq}
Was the $\tau$ search grid extended beyond 2.0 for the $\tau = 2.4$ set? If
it was not, every fit in that set is boundary-biased in $\tau$ and part of
the AICc collapse there is a grid-capping artefact. The trend is established
within bounds at $\tau = 1.8$ and so survives either way, but the $\tau=2.4$
magnitudes depend on the answer. If the grid was capped, describe that set
as a stress test beyond the operational bound.
\end{openq}

\begin{openq}
The \citet{DeSilva:2019aa} comparison is not like for like. Their figure is
for prediction of the later curve from data prior to $-2$~LU; ours is
estimation from the whole curve. Either state the difference explicitly or
run the prediction version, truncating the simulated curves at $-2$~LU and
recomputing RIT error, which would make the comparison exact rather than
indicative.
\end{openq}

\begin{towrite}
Remaining limitations: simulation uses the same generative model as the fit,
so model misspecification is untested; noise is additive Gaussian on
threshold, and real thresholds are quantised by the staircase; no clinical
data appears in this paper.
\end{towrite}

% ================= 9 =================
\section{Conclusion}
\begin{towrite}
Three sentences. Separability makes the fit deterministic and cheap; the
validation quantifies exactly what that buys and what it does not; the rod
slope is a usable clinical parameter with a stated detection threshold.
\end{towrite}

% ================= back matter =================
\section*{Data and code availability}
\begin{towrite}
State the analysis scripts by name (\texttt{dc\_findInitialParametersProfiled.R},
\texttt{make\_result\_row.R}, \texttt{dc\_analysis.R},
\texttt{dc\_run\_experiment\_driver.R}, \texttt{dc\_envCap.R},
\texttt{dc\_platform\_compare.R}, \texttt{dc\_sparse\_study.R},
\texttt{dc\_rit\_error\_figure.R}) and the commit against which the reported
results were produced. Decide what is released and what is withheld given
the pending disclosure.
\end{towrite}

\section*{Competing interests}
\begin{towrite}
JMFK is a director and shareholder of MuMac Ltd and an inventor on
\citet{Kelly:2020aa}. Declare the pending disclosure.
\end{towrite}

\bibliography{varpro}

\end{document}
